% ════════════════════════════════════════════════════════════
\section{CAPÍTULO I: PLANTEAMIENTO DEL PROBLEMA}
% ════════════════════════════════════════════════════════════

\subsection{Planteamiento de problema}

El desarrollo de software moderno depende casi universalmente de paquetes de terceros.
Cuando un desarrollador instala una librería desde el \textit{Python Package Index}
(PyPI), no instala únicamente ese paquete: sino que instala un árbol de dependencias que puede
involucrar decenas o cientos de paquetes adicionales, cada uno mantenido por distintos
equipos con distintos niveles de actividad y seguridad. Este fenómeno genera lo que
la literatura de ingeniería de software denomina \textit{supply chain risk}: el riesgo
de que un fallo en un componente aparentemente secundario se propague en cascada y
comprometa a miles de proyectos que dependen de él \autocite{zimmermann2019small}.

El problema central es de naturaleza estructural y probabilística: dada una red de
$n$ paquetes interconectados mediante relaciones de dependencia, \textbf{¿cómo
identificar qué paquetes representan un riesgo sistémico desproporcionado para el
ecosistema, antes de que una vulnerabilidad sea registrada?}

Casos reales ilustran la gravedad del problema. En 2016, la eliminación del paquete
\texttt{left-pad} (11 líneas de código) de npm provocó que React, Babel y miles de
proyectos de producción dejaran de compilar en todo el mundo durante horas
\autocite{decan2019impact}. En 2021, la vulnerabilidad en \texttt{log4j} expuso
servidores de Amazon, Apple y Cloudflare porque nadie tenía un mapa claro del
alcance transitivo de esa dependencia \autocite{cve-2021-44228}. En ambos casos, el problema no era la
ausencia de herramientas de escaneo, sino la ausencia de una \textbf{métrica
global de criticidad sistémica} que hubiera permitido priorizar auditorías antes
del incidente.

Las herramientas actuales (Dependabot, Snyk, \texttt{pip-audit}) analizan proyectos
individualmente y de forma reactiva: actúan cuando ya existe un CVE registrado.
Ninguna modela el ecosistema completo como una red y ninguna estima la propagación
transitiva del riesgo sobre su estructura.

\subsection{Justificación del problema}

\subsubsection{Justificación teórica}
El grafo de dependencias de PyPI es un sistema complejo cuya estructura determina
la dinámica de propagación de fallos. Modelarlo como una Cadena de Markov en tiempo
discreto permite aplicar directamente los teoremas de Perron–Frobenius, la propiedad
de ergodicidad y la convergencia del método de la potencia a un problema de
ingeniería de software relevante y abierto. El PageRank, como distribución
estacionaria de dicha cadena, captura tanto la centralidad directa como la
centralidad transitiva de cada paquete, lo que lo convierte en un candidato
natural para medir riesgo sistémico en redes de dependencias.

\subsubsection{Justificación práctica}
La seguridad de la cadena de suministro de software (\textit{software supply chain
security}) es una de las áreas de mayor actividad en la industria. El gobierno de
Estados Unidos emitió en 2021 una orden ejecutiva que exige a los proveedores de
software generar \textit{Software Bills of Materials} (SBOMs) para identificar
dependencias \autocite{kikas2017structure}. Una métrica de PageRank sobre PyPI
ofrecería a auditores, mantenedores y organizaciones de seguridad una herramienta
proactiva para priorizar esfuerzos antes de que ocurran incidentes.

\subsubsection{Justificación metodológica}
El trabajo integra tres ejes del curso en un pipeline coherente:
\begin{enumerate}
    \item \textbf{Análisis exploratorio de grafos}: caracterización topológica del
          grafo de dependencias (distribución de grado, componentes, densidad).
    \item \textbf{Cadenas de Markov}: construcción de la matriz de transición,
          verificación de ergodicidad y cálculo de la distribución estacionaria
          (PageRank).
    \item \textbf{Análisis espacial}: detección de comunidades y visualización
          del grafo como un espacio de clusters temáticos del ecosistema.
\end{enumerate}

\subsection{Determinación, delimitación del problema}

\begin{description}[style=nextline]
    \item[Temática]
        Aplicación de Cadenas de Markov en tiempo discreto y PageRank para la
        identificación de paquetes sistémicamente críticos en PyPI. No se abordan
        extensiones como PageRank personalizado, análisis de vulnerabilidades
        específicas (CVEs), ni métricas de calidad de código de las dependencias.

    \item[Datos]
        Se utilizará el dataset de dependencias de PyPI disponible en
        \textit{libraries.io} \autocite{librariesio}, que registra las relaciones
        de dependencia entre paquetes publicados en PyPI. El análisis se
        restringirá a la componente gigante del grafo para garantizar la
        unicidad de la distribución estacionaria antes de aplicar la corrección
        de \textit{teleportation}.

    \item[Temporal]
        Snapshot estático del ecosistema (sin modelado de evolución temporal
        del grafo de dependencias).

    \item[Computacional]
        Implementación en Python 3.11 con las bibliotecas \texttt{NumPy},
        \texttt{SciPy} (álgebra lineal dispersa), \texttt{NetworkX} y
        \texttt{python-igraph} para análisis de comunidades.
\end{description}

\subsection{Interrogantes del problema}

\begin{quote}
\textbf{Problema general:}
¿En qué medida el PageRank calculado sobre el grafo de dependencias de PyPI,
fundamentado en la distribución estacionaria de una Cadena de Markov, constituye
un mejor predictor del impacto sistémico de un fallo que las métricas locales
disponibles actualmente (número de descargas, número de dependientes directos,
presencia de CVEs)?
\end{quote}

\begin{quote}
\textbf{Problemas específicos:}
\begin{enumerate}[label=\textbf{PE\arabic*.}]
    \item ¿Cuál es la estructura topológica del grafo de dependencias de PyPI y bajo
          qué condiciones la Cadena de Markov asociada garantiza una distribución
          estacionaria única?
    \item ¿Qué paquetes obtienen mayor PageRank en el ecosistema PyPI y en qué medida
          este ranking difiere del ordenamiento por métricas locales (in-degree,
          número de descargas)?
    \item ¿Cómo se distribuyen los paquetes de alto PageRank entre los distintos
          subcomunidades (clusters) del ecosistema, y qué estructura espacial
          presenta la red?
    \item ¿Cuántos proyectos activos de PyPI dependen transitivamente de los
          10 paquetes con mayor PageRank, y qué proporción representa respecto
          al total del ecosistema?
\end{enumerate}
\end{quote}

\subsection{Objetivos}

\subsubsection{Objetivo General}
Identificar los paquetes sistémicamente críticos del ecosistema PyPI mediante el
cálculo del vector PageRank sobre su grafo de dependencias, modelado como la
distribución estacionaria de una Cadena de Markov, y comparar este ranking con
métricas locales de centralidad para evaluar su capacidad predictiva del riesgo
sistémico.

\subsubsection{Objetivos Específicos}
\begin{enumerate}[label=\textbf{OE\arabic*.}]
    \item Construir el grafo dirigido de dependencias de PyPI a partir del dataset
          de \textit{libraries.io} y caracterizar su topología mediante métricas
          exploratorias: distribución de grado, densidad, coeficiente de
          agrupamiento y componentes conexas.

    \item Formalizar la Cadena de Markov sobre el grafo de dependencias mediante
          la construcción de la matriz de transición $\mathbf{H}$, la corrección
          de nodos sumidero y la matriz de Google $\mathbf{G}(d)$, verificando
          las condiciones de ergodicidad mediante el Teorema de Perron–Frobenius.

    \item Calcular el vector PageRank mediante el método de la potencia iterativo
          para $d \in \{0.75, 0.85, 0.90, 0.95\}$ y analizar la sensibilidad del
          ranking Top-20 al factor de amortiguamiento.

    \item Detectar comunidades en el grafo de dependencias mediante el algoritmo
          de Louvain y analizar la distribución espacial de los paquetes de mayor
          PageRank entre los clusters identificados.

    \item Comparar el ranking por PageRank con los rankings por in-degree,
          número de descargas y número de dependientes directos, cuantificando
          las diferencias mediante correlación de Spearman $\rho_s$.
\end{enumerate}

\subsection{Variables}

\subsubsection{Variables Dependientes}
\begin{center}
\begin{tabular}{p{3.5cm} p{9cm}}
\toprule
\textbf{Nombre} & \textbf{Vector PageRank} $\mathbf{r} = (r_1, r_2, \ldots, r_n)^\top$ \\
\midrule
Definición conceptual &
    Importancia sistémica global de cada paquete, definida como la probabilidad
    de que un agente aleatorio que navega el grafo de dependencias se encuentre
    en dicho paquete en el estado estacionario. \\[6pt]
Definición operacional &
    Componente $i$-ésima del vector propio dominante (autovalor $\lambda_1 = 1$)
    de la matriz de Google $\mathbf{G}(d)$:
    \[
        \mathbf{r}^* = \lim_{k \to \infty} \mathbf{G}(d)^k \, \mathbf{r}^{(0)},
        \quad \mathbf{r}^{(0)} = \tfrac{1}{n}\mathbf{1}
    \]  \\[6pt]
Tipo & Cuantitativa continua (probabilidad: $r_i \in [0,1]$, $\sum_i r_i = 1$) \\[6pt]
Escala & Razón \\[6pt]
Indicadores &
    \begin{itemize}[nosep]
        \item Rango ordinal del paquete (posición en el ranking global)
        \item Correlación de Spearman $\rho_s$ vs.\ métricas locales
        \item Proporción del ecosistema cubierta por los Top-$k$ paquetes
    \end{itemize} \\
\bottomrule
\end{tabular}
\end{center}

\subsubsection{Variables Independientes}
\begin{center}
\begin{tabular}{p{3.5cm} p{3.5cm} p{2.5cm} p{3cm}}
\toprule
\textbf{Variable} & \textbf{Definición operacional} & \textbf{Tipo} & \textbf{Rango / Valores} \\
\midrule
Factor de amortiguamiento $d$ &
    Probabilidad de que el agente siga una dependencia en lugar de teleportarse;
    controla la mezcla entre $\mathbf{H}$ y la distribución uniforme.
    $\mathbf{G}(d) = d\mathbf{A} + (1-d)\mathbf{E}/n$ &
    Cuantitativa continua &
    $\{0.75,\, 0.85,\, 0.90,\, 0.95\}$ \\[10pt]
In-degree del paquete &
    Número de paquetes que lo declaran como dependencia directa. &
    Cuantitativa discreta &
    $k^{\text{in}} \in \mathbb{N}_0$ \\[10pt]
Número de descargas &
    Total de descargas en PyPI (últimos 30 días, fuente BigQuery). &
    Cuantitativa discreta &
    $\geq 0$ \\[10pt]
Estructura del grafo $G=(V,E)$ &
    Número de paquetes $|V|=n$ y dependencias $|E|=m$; densidad $\delta = m/n^2$. &
    Cuantitativa (fija) &
    Dataset libraries.io \\
\bottomrule
\end{tabular}
\end{center}

\subsubsection{Variable Interviniente}
\begin{description}
    \item[Nodos sumidero (\textit{dangling nodes})] Paquetes sin dependencias
        salientes (\texttt{out-degree}$=0$), es decir, paquetes que no requieren
        ninguna otra librería. Representan un sumidero en la cadena que rompe la
        condición de estocasticidad de $\mathbf{H}$. Se corrigen reemplazando
        su columna por $\mathbf{1}/n$ para obtener la matriz estocástica
        $\mathbf{A}$.
\end{description}

\subsection{Hipótesis}

\subsubsection{Hipótesis General}
\begin{quote}
El PageRank calculado sobre el grafo de dependencias de PyPI, como distribución
estacionaria de una Cadena de Markov ergódica, identifica un conjunto de paquetes
cuyo impacto transitivo sobre el ecosistema es significativamente mayor al
predicho por métricas locales (in-degree, número de descargas), evidenciado por
una correlación de Spearman $\rho_s < 0.80$ entre el ranking por PageRank y
el ranking por in-degree normalizado.
\end{quote}

\subsubsection{Hipótesis Específicas}
\begin{enumerate}[label=\textbf{H\arabic*.}]
    \item El grafo de dependencias de PyPI presenta una distribución de grado con
          cola pesada (ley de potencias), lo que implica la existencia de un
          reducido conjunto de paquetes que concentran la mayor parte de las
          dependencias del ecosistema.
    \item Los 10 paquetes con mayor PageRank cubren, de forma transitiva, más
          del 60\% de los proyectos activos en PyPI, evidenciando un riesgo
          sistémico altamente concentrado.
    \item Existen paquetes con PageRank significativamente superior a su in-degree
          normalizado, lo que indica que reciben masa de probabilidad de forma
          transitiva a través de dependencias de alta centralidad.
    \item Los paquetes de mayor PageRank se concentran en un número reducido de
          comunidades del grafo (clusters del ecosistema data science, web y
          criptografía), revelando una estructura espacial no aleatoria del riesgo.
\end{enumerate}
